%!TeX root=../thesis.tex

\chapter{Research Instruments}
\label{apx:instruments}

This appendix reproduces the instruments and communications used during data
collection. Section~\ref{sec:listening-guide} contains the open-ended active
listening guide used in the end-of-period sessions.
Section~\ref{sec:informed-consent} reproduces the informed consent framework
communicated to participants. Section~\ref{sec:training-comm} reproduces the
training communication sent to both teams at the start of the observation period.

% ─────────────────────────────────────────────────────
\section{Open-Ended Active Listening Guide}
\label{sec:listening-guide}

\textbf{Note on instrument design.} The study did not employ a fixed
questionnaire. As described in Section~\ref{subsec:active-listening}, data collection followed an
open-ended active listening format consistent with the semi-structured
interview approach recommended by \citet{robson2016} for flexible designs.
The guide below provided a thematic scaffold for conversations; the actual
order and phrasing of questions varied with the natural flow of each session.
Sessions were conducted individually with the senior technical leader of each
team at the end of the observation period.

\subsection*{Thematic Areas and Guiding Questions}

\textbf{Area 1 --- Overall experience of Hypo-Stage.}
\begin{itemize}
  \item Can you walk me through the moments when you actually opened Hypo-Stage
    during the observation period?
  \item What made you decide to register a hypothesis at those particular
    moments?
  \item Were there situations where you thought about registering a hypothesis
    but ended up not doing so? What stopped you?
\end{itemize}

\textbf{Area 2 --- Perceived usefulness of the technique and the tool.}
\begin{itemize}
  \item Has using Hypo-Stage changed the way you think about or discuss
    architectural decisions with your team?
  \item Are there specific decisions you handled differently because you had a
    hypothesis registered?
  \item What aspect of the tool did you find most useful in practice?
  \item What aspect of the tool was most frustrating or felt like friction?
\end{itemize}

\textbf{Area 3 --- Learning curve and technique understanding.}
\begin{itemize}
  \item How confident do you feel formulating a hypothesis statement that is
    falsifiable? Did that feel natural or forced?
  \item How did your team members react to the concept of an architectural
    hypothesis? Do you think the concept would be accessible to a mid-level
    engineer without facilitation?
  \item What would have made the onboarding experience easier?
\end{itemize}

\textbf{Area 4 --- Workflow integration.}
\begin{itemize}
  \item How well did Hypo-Stage fit into your existing sprint ceremonies and
    day-to-day tooling?
  \item Did the Backstage integration help or was it irrelevant to adoption?
  \item If you were to continue using the tool after this study, what would
    need to change --- either in the tool or in your team's process --- for it to
    be sustainable?
\end{itemize}

\textbf{Area 5 --- Future outlook.}
\begin{itemize}
  \item If you had to name one architectural uncertainty that you are currently
    managing tacitly that you \emph{wish} you had registered in Hypo-Stage,
    what would it be?
  \item Would you recommend Hypo-Stage to another team in the organization?
    Under what conditions?
\end{itemize}

% ─────────────────────────────────────────────────────
\section{Informed Consent Framework}
\label{sec:informed-consent}

No formal written consent form was required for this study, as the research
involved no sensitive personal data and participation was entirely voluntary
(Section~\ref{sec:ethical-considerations}). The following points were communicated verbally and in writing
to all participants prior to the observation period and were confirmed by
engineering leadership in the organization:

\begin{enumerate}
  \item \textbf{Voluntary participation.} All engineers were free to use
    Hypo-Stage as much or as little as they chose. Participation in the
    active listening sessions was optional and had no bearing on the
    researcher's organizational role.
  \item \textbf{Anonymization.} The name of the organization, the names of
    all participants, and the names of internal software products and
    components would be anonymized in all research outputs. Backstage entity
    references in hypothesis data would be replaced with generic labels
    (see Appendix~\ref{app:hypotheses-data}, Table~\ref{tab:entity-legend}).
  \item \textbf{Data use.} Observations, artifact records, and session notes
    would be used exclusively for academic research purposes and would not be
    shared with organizational management in any individually identifiable form.
  \item \textbf{Right to withdraw.} Any participant could withdraw previously
    provided data or discontinue participation at any point without consequence.
  \item \textbf{Research knowledge and consent of leadership.} The study
    was conducted with the knowledge and informal consent of the relevant
    engineering leadership in the organization, who confirmed that participation
    was compatible with the teams' work obligations.
\end{enumerate}

% ─────────────────────────────────────────────────────
\section{Training Communication}
\label{sec:training-comm}
\label{apx:training-communication}

The following text reproduces, in anonymized and lightly edited form, the
written communication sent to both teams prior to the training session that
inaugurated the observation period. The purpose of this communication was to
set expectations, introduce the research context, and invite voluntary
participation.

\begin{quote}
\textbf{Subject: Architectural Hypothesis Management with Hypo-Stage ---
Invitation to Participate}

Hi Team,

I am reaching out to invite you to participate in a research initiative I am
conducting as part of my Master's degree at IME-USP.

\textbf{What is it about?}
We will be exploring a technique called Architectural Hypothesis Engineering
(ArchHypo) and a supporting tool called Hypo-Stage, which I have integrated
into our Backstage portal. The idea is simple: instead of leaving
architectural uncertainties as tacit knowledge held by a few engineers, we
write them down as structured hypotheses, assess their urgency, and track
the actions we are taking to resolve them.

\textbf{What does participation involve?}
I will run a short training session (approximately 60--90 minutes) to introduce
the technique and the tool. After that, you are invited to use Hypo-Stage
whenever you encounter architectural uncertainties in your normal work over
the next four weeks. At the end of the period, I would like to have a brief
conversation (20--30 minutes) with interested team members about their
experience.

\textbf{Is it mandatory?}
No. Participation is entirely voluntary. Using the tool as much or as little
as you find valuable is up to you.

\textbf{What happens to the data?}
All records will be anonymized before being reported in my thesis. The names
of our teams, our products, and our systems will not appear in any research
output. Observations are used only for academic research.

Please reply if you have any questions, or let me know if you would like to
discuss the initiative before the training session.

Looking forward to working with you on this!

[Researcher --- role and name anonymized]
\end{quote}

\paragraph{Editorial note on observation length.}
The invitation above references a four-week horizon. In practice, teams continued to use Hypo-Stage voluntarily beyond that window. Hypothesis registration timestamps in Appendix~\ref{app:hypotheses-data} span 11~February~2026 to 12~April~2026 (61~calendar days, approximately nine weeks), corresponding to five consecutive two-week sprint periods from the earliest to the latest registration in the organization's cadence.
